\section{Actual paired lifts and finite cleanup}
\label{sec:paired-lifts}

We now construct the two polynomials used in the Type-I pairing. Work
in the literal unramified quadratic setting, with uniformizer \(w^2=p\),
\(v_p(p)=1\), \(k=m+1\), and \(Q=p^{G'}\). Assume the actual
\(\mathrm{Std}\) hypotheses, \(C_{E1},C_{E2}\ge17\), \(C_m\ge17\).
The kept-block, terminal-degree and order bounds are those of SW12;
the coarse-degree alternative is SW10. The endpoint height below is
\(t_0\); its subsequent comparison with \(t_I\) is separate.

\paragraph{RL1--RL3: correction of the actual segment form.}
For a block \(b\), the physical monomial substitution is
\[
 I_b(\alpha)=\rho_b+\sum_v P_v\alpha_v e_{\iota(v)}.
\]
The map \(\iota\) is injective, \(P_v=p^{g_v}\) at odd positions and
\(P_v=1\) at the terminal position. Balanced degrees imply that the
total auxiliary degree \(c_\ell\) at each physical letter, and the total
degree at each radix \(P_v\), are fixed. Put
\[
 C_b=\prod_\ell c_\ell!,\quad N_b=C_b\,\mathrm{sp}_b,\quad
 v_b=v_p(N_b),\qquad
 L_j=\sum_\ell\lfloor\log_p d_\ell\rfloor,\quad
 z_j=\frac{p^{kL_j}}{\prod_\ell d_\ell!}.
\]
Here \(\mathrm{sp}_b\) is the source's integer \(\mathrm{specConst}\);
zero degrees contribute zero to \(L_j\).
DT's inverse factorial-unit approximation supplies
\[
 g_b(\alpha)=\sum_{|\nu|\le4c}q_\nu\binom{\alpha}{\nu},\qquad
 I_b(\alpha)!\,g_b(\alpha)
   =\Phi_b\alpha!(1+p^ch_\alpha),\quad q_\nu,h_\alpha\in O.
 \tag{RL.1}
\]
Only odd variables occur in \(\nu\). Explicitly,
\(\Phi_b=R_b\prod_v\rho_v!(P_v!)^{\alpha^0_v}\), where \(R_b\)
contains the fixed factorials outside \(\iota\), and \(\alpha^0\)
puts each normal-letter degree in its first coordinate.

The scalar
\[
 \kappa_b=\frac{z_jp^{v_b}\Phi_b}{\mathrm{sp}_b}
 \quad\text{satisfies}\quad
 0\le v_p(\kappa_b)\le(k+1)L_j.                  \tag{RL.2}
\]
Indeed \(C_b/\alpha^0!\) is a product of binomial coefficients, whereas
\(\prod d_\ell!/I_b(\alpha^0)!\) is a product of \(k\)-part multinomials.
Their valuations are at most \(L_j\) and \((k-1)L_j\), respectively.
Legendre's formula makes \(v_p(\kappa_b)\) equal to \(kL_j\) plus the
first valuation minus the second.

Replace each binomial in \(g_b\) by the diagonal operator
\(x^\nu\partial^{[\nu]}\), obtaining \(G_b\), and set
\[
 s_b(u)=w^{v_b}G_bu,\qquad s'_b(u)=w^{v_b}u.       \tag{RL.3}
\]
The weighted coefficient condition is
\(v_p(u_\alpha)+\tfrac12v_p(\alpha!\mathrm{sp}_b)\ge0\).
Since \(\alpha!\mathrm{sp}_b\mid N_b\), both outputs are integral.
These operators reduce the SSZ margin by at most \(4c\): expand the
coefficient of \(t^n\) under \(x\mapsto(1+t)x\), and in each tied factor
use \(ZS=E-q\), where \(E=ZS+q\) and \(q=\sum_aW_aE_a\).
Each changed coordinate or separated pair costs at most one prefix
weight. All binomial coefficients are integral; SW1 supplies saturation.
Thus the maps land in \(\mathrm{SSZ}(m_0(13))\), because
\[
 m_0(C_m)-4c\ge m_0(13),\qquad
 m_0(C)=CkN_1Qc+12c .
\]
Equation (RL.1), summed coefficientwise, gives
\[
 z_j\operatorname{Fock}(\operatorname{io}_bs_bu,
                        \operatorname{io}_bs'_bv)
 =\kappa_b\bigl(\mathrm{sp}_b\operatorname{Fock}(u,v)+w^c\eta\bigr).
 \tag{RL.4}
\]
Distinct blocks are exactly orthogonal: reduce physical exponents modulo
\(P_i\), recover the early coordinate-zero degrees, then recover the
remaining block orders from the balanced degrees. Segments use disjoint
physical letters. Moreover \(z_{\rm Ten}\alpha!\in O\) for every physical
monomial of the fixed degree, by the same multinomial valuation bound.
Consequently (RL.4) factors over segments and is unchanged modulo \(w^c\)
when either symbol is changed modulo \(w^c\).

\paragraph{RL4--RL7: actual raw polynomials and their tail spaces.}
Put \(A=13kN_1Q\), \(\mu=mm+nA\), \(E=(c-n)A\), where
\(mm\in\{0,12c\}\). The numerical wrappers admit explicit short estimates.
With \(\epsilon=N_1\mathrm{La}/k\), the early coarse alternative and
\(\mathrm{Std}\) give
\[
 \mu+N_1E\le4\epsilon A_0,\quad
 \sum_{i<\sigma}\mathrm{sg}_i\le2\epsilon A_0+N_1E,\quad
 \sum_{\ell\le i}w_\ell\le A_i(1+\epsilon),\quad
 A_i(1+\delta)<A_0 .
\]
Since \(40\epsilon\le\delta<1/6\), these prove the literal
floor inequalities \(\mathrm{ColumnNumericsAt}\).
IP's joined one-plane/normal paths therefore give actual Frobenius
column polynomials, of letter degree at most \(Q\).

Multiply these columns by the Frobenius minors and contractions
\[
 K_{W,a}=X_{\mathrm{odd},+,a}X_{\mathrm{top},+,0}^{p^d}
          -X_{\mathrm{odd},+,0}X_{\mathrm{top},+,a}^{p^d},
 \qquad K_S=\sum_aX_{\mathrm{top},+,a}^{p^d}X_{\mathrm{odd},-,a},
\]
with the recorded digit and level multiplicities. Include
\[
 \left[X_{\mathrm{top}_{j_0},+,0}
 \prod_{a>0}\left\{\prod_jX_{\mathrm{top}_j,+,0}^{p^{G_j}-1}
                    -\prod_jX_{\mathrm{top}_j,+,a}^{p^{G_j}-1}
             \right\}\right]^c.                         \tag{RL.5}
\]
This localizer vanishes modulo \(w\) at every other line, by cyclic
Frobenius telescoping. A residual monomial makes the own leading term
exactly the prescribed symbol. Its exponents are nonnegative because
\[
 \mathrm{usedX}\le a-E_0+(N_1+c)kQ,\qquad
 E_0\ge\mathrm{localOrderCap}+cA+(N_1+c)kQ .
\]

For tails removing at most \(e\) factors, divide each new exponent by
\(P_i\), keeping both quotient and remainder. The remainder defines
the new digits. The quotients give a new block, and
\(Z^zS^y=Z^{z-y}(E-q)^y\) gives its exact integral expansion.
The required \(z\ge y\) follows from
\[
 \delta A_i/4\ge c(A+2Q+3k)+4.
\]
Deleting columns costs \(N_1e\) in the prefix bound; the new monomial
and tied factors cost at most \((N_1+1)(k+Qe)\). The cap decreases by
at most \(N_1(3k+2Qe)\), and
\[
 N_1e+(N_1+1)(k+Qe)+N_1(3k+2Qe)\le eA.          \tag{RL.6}
\]
Taking floors preserves this integer inequality. This proves the
actual \(\mathrm{NumCoreAt}\) and tail-space membership, including
the natural-number subtractions.

\paragraph{RL8--RL11: finite cleanup and the genuine paired endpoint.}
Write the product as \(f_R\prod_af_a^{m_a}\), with exact rigid
chart factor and own leading terms \(w^{o_a}\sigma_a\).
For \(t_a=w^{-o_a}\Phi f_a-\sigma_a\),
\[
 v_w([X^M]t_a)\ge\max(1,\operatorname{nd}(M)-o_a).
\]
Expansion and (RL.6) give a raw lift of height \(o\) satisfying
\[
 \Phi(\mathrm{raw})=\sigma+\sum_{h=1}^{n-1}w^h\tau_h\pmod{w^n},
 \quad \tau_h\in\mathrm{SSg}(mm+(n-h)A,(c-n+h)A,o+h).
\]
Inductively lift \(\tau_h\) to precision \(n-h\) and subtract
\(\sum_hw^hF_h\). Multiplication by \(w^h\) cancels its additional
height \(h\); all other-edge conditions remain divisible by \(w^c\).
At \(n=c\) this produces actual cleaned lifts of height
\(o\le\lfloor t_0(C_0,C_t)\rfloor\), \(C_t\ge30\).

For the edge basis \(b\), write
\[
 b'=b^{-T}\operatorname{diag}(p,1,\ldots,1),\qquad
 D_w=\operatorname{diag}(w^{-1},1,\ldots,1).
\]
The reversed dual midpoint satisfies \(b'D_w=(bD_w)^{-T}\).
The actual flat maps and normalized Fock map
\(\Xi\) therefore satisfy
\[
 \mathrm{beHH}(\Theta_bF,\Xi\Theta_{b'}F')
          =z_{\rm Ten}\operatorname{Fock}(\Phi_1F,\Phi_1F').
\]
Normalized integrality sends the dual cleaned lift into
\(\mathrm{CCe}'\). The resulting actual vectors satisfy
\[
 \begin{gathered}
 F_\beta(u)\in\mathrm{CCe},\quad G_\beta(u)\in\mathrm{CCe}',\\
 \mathrm{beHH}(F_\beta(u),G_{\beta'}(u'))
 =\mathbf1_{\beta=\beta'}\kappa_\beta
       \prod_j\mathrm{segForm}_j(u_j,u'_j)+w^c\nu ,\\
 v_p(\kappa_\beta)\le4\kappa_0.
 \end{gathered}                                  \tag{RL.7}
\]
Nonretained block components are zero. This proves the literal
\(\mathrm{PairLiftAt}\) conclusion; the separate digit-class rank
and subsequent global parameter choices are not inferred from it.
